\chapter{Introduction}
\label{chap:introduction}

\section{Background}
\label{section:background}

In everyday life, people rely on multiple sensory inputs to understand and interact with their surroundings. 
Among these, hearing plays an important role in recognizing environmental events such as alarms, 
approaching vehicles, and nearby activities. These auditory cues provide important information 
that helps individuals interpret what is happening around them.

However, individuals who are deaf or hard of hearing are unable to perceive these sounds, 
which can make it more difficult to access this information. As a result, 
they often depend more on visual cues or external assistance to understand their environment.

Assistive technologies such as hearing aids and cochlear implants have been developed to address this limitation. 
While these technologies can be effective for some users, 
they are not suitable for everyone and may not fully capture all types of environmental sounds. 
This highlights the need for alternative approaches that do not rely solely on restoring hearing.

This project proposes an application that represents environmental sounds in a visual form. 
It is not intended to function as a primary or continuous tool for safety-critical awareness, 
as relying solely on a visual interface requires constant attention and may not be appropriate in situations that demand immediate reactions, 
such as navigating traffic. For this reason, the application is designed as a supplementary tool rather than a replacement for existing assistive technologies. 
It enables users to observe and explore sound information by capturing audio input and presenting it visually, 
including properties such as sound type, direction, and distance through a spatial, map-like interface. 
The primary focus is to enhance quality of life in non-critical situations by providing an additional way to understand surrounding sounds.

While the current application is designed as a supplementary tool, it also serves as a foundation for future development. 
The underlying concept of visualizing sound could be further adapted into more integrated and responsive systems, 
such as wearable devices or embedded assistive technologies, where information can be delivered more seamlessly. 
With such advancements, this approach has the potential to evolve beyond a supportive application and contribute to more comprehensive assistive solutions.

\section{Problem Statement}
\label{section:problem-statement}

Individuals who are deaf or hard of hearing face difficulties in perceiving environmental sounds that provide important contextual information about their surroundings. 
While existing assistive technologies such as hearing aids and cochlear implants can enhance auditory perception, 
they are not suitable or effective for all users and may not capture all types of environmental sounds.

As a result, there is a lack of accessible solutions that allow these individuals to understand surrounding audio information in an alternative and intuitive way. 
In particular, there is limited support for representing sound characteristics, such as type, direction, and distance, 
in a form that can be easily interpreted without relying on auditory input.

Additionally, relying solely on visual tools for real-time awareness in safety-critical situations may not be appropriate, 
as such approaches require continuous attention and may introduce risks in dynamic environments. 
Therefore, there is a need for a solution that provides meaningful sound awareness while being suitable for non-critical, 
everyday contexts.

This project addresses the problem by exploring how environmental sound can be transformed into visual representations 
to support better understanding and improve quality of life for users who cannot depend on traditional hearing-based technologies.

\section{Solution Overview}
\label{section:solution-overview}

The proposed solution is a software application that transforms environmental audio into visual information to assist users in understanding surrounding sounds. 
The system captures real-time audio input and processes it to identify key characteristics such as sound type, direction, 
and relative distance. This information is then presented through a spatial, map-like interface, 
allowing users to intuitively interpret where sounds originate and what they represent.

The application is designed to provide a supplementary method of interaction with the environment, 
particularly for individuals who are deaf or hard of hearing and may not benefit fully from traditional assistive technologies. 
By converting sound into visual cues, the system offers an alternative way to access information that is typically conveyed through hearing.

The primary focus of the solution is to enhance quality of life in non-critical situations, 
where users may wish to explore or better understand their surroundings. 
The system is not intended to replace existing assistive tools or serve as a primary solution for safety-critical awareness. 
Instead, it complements existing methods by providing an additional layer of information in a user-friendly format.

In addition, the application is designed with extensibility in mind. 
The core concept of visualizing sound can be further developed and adapted into other platforms, 
such as wearable devices or integrated assistive systems, 
enabling more seamless and context-aware delivery of information in the future.

\subsection{Features}
\label{subsection:features}

\begin{enumerate}[leftmargin=80pt]
    \item Real-time Sound Detection: Captures environmental audio input continuously from the device microphone.
    \item Sound Classification: Identifies and categorizes detected sounds (e.g., speech, vehicle, alarm, ambient noise).
    \item Direction Detection: Estimates the direction from which a sound originates relative to the user (e.g., east, north, east-west).
    \item Distance Level Estimation: Provides an approximate indication of the sound source proximity using relative levels (e.g., near, far, very far).
    \item Visual Map Interface: Displays detected sounds on a spatial, map-like interface for intuitive understanding.
    \item Sound Visualization Indicators: Uses visual elements to represent different sound types and properties (e.g., icons, colors, markers).
    \item User Interaction: Allows users to view, explore, and interpret sound information through the interface.
    \item Non-Critical Usage Mode: Designed for use in non-safety-critical situations to enhance quality of life rather than replace primary assistive tools.
    \item Extensibility Support: Designed to allow future integration with wearable or assistive technologies.
\end{enumerate}

\section{Target User}
\label{section:target-user}

The primary target users of this application are individuals who are deaf or hard of hearing, 
particularly those who are unable to fully benefit from traditional assistive technologies such as hearing aids. 
These users may seek alternative ways to understand environmental sounds in their surroundings through visual representations.

In terms of demographics, the application is intended for a broad age range, including teenagers and adults, 
with basic familiarity with smartphones or similar devices. 
The design assumes general user accessibility rather than targeting a highly specialized professional group.

Regarding skill level, the application is designed for users with low to moderate technical proficiency. 
The interface emphasizes simplicity and intuitive interaction, 
minimizing the need for extensive training or prior experience with similar systems.

The application is not limited to a specific industry or domain, 
but is intended for general everyday use in non-critical situations. 

It may also be useful for individuals without hearing impairments who wish to explore environmental sound information through visual representation. 
For example, users may use the application to identify different sound sources in environments such as busy streets, public spaces, or natural settings. 
The application can also be beneficial in situations where audio perception is limited or inconvenient, 
such as in very noisy environments or quiet zones. Additionally, it may serve as an educational or exploratory tool, 
helping users better understand how sound behaves in different contexts, including its direction and relative distance.

Overall, the system is designed to be accessible, easy to use, 
and adaptable to the needs of users who require alternative methods for perceiving and understanding sound.

\section{Terminology}
\label{section:terminology}

\begin{enumerate}[leftmargin=80pt]
    \item Environmental Sound: Any sound occurring in the surrounding environment, such as speech, vehicles, alarms, or background noise.
    \item Sound Classification: The process of identifying and categorizing detected audio into predefined types (e.g., speech, traffic, alarm).
    \item Sound Direction: The relative position from which a sound originates with respect to the user (e.g., east, north, east-west).
    \item Distance Level: A qualitative indication of how close or far a sound source is from the user (e.g., near, far, very far).
    \item Sound Visualization: The representation of audio information in a visual format, such as icons, markers, or spatial elements.
    \item Spatial Interface: A visual layout that represents the relative position of elements in space around the user, often using a map-like or abstract view.
    \item Non-Critical Usage: Situations where the system is used for general awareness or exploration rather than safety-critical decision-making.
    \item Supplementary Tool: A system designed to support or enhance existing methods, rather than replace them.
    \item Real-time Processing: The ability of the system to analyze and present data immediately or with minimal delay.
\end{enumerate}